% GENERATED FILE. Edit canonical lessons, then run npm run generate:books.
% canonical-source-hash: fnv1a64:b15bd550928701ec
% canonical-lessons: MR-C286-jar-tar, MR-C286-padlyas, MR-R286-if-and-otherwise

\chapter{If, and Otherwise}
\label{ch:mr-if-and-otherwise}
\hlchaptermodality{286}

\begin{chapteropening}
\textbf{By the end of this chapter:} \emph{I can say what will happen if something does (jar … tar …), read the formal -lyās, and warn with nāhītar.}

Complete MR-R286-if-and-otherwise: set a condition and its result, and say what happens otherwise.
\end{chapteropening}

\section[\texorpdfstring{\mr{जर} … \mr{तर} … (jar … tar …)}{जर … तर … (jar … tar …)}]{\mr{जर} … \mr{तर} … (jar … tar …) — if}

\label{lesson:MR-C286-jar-tar}



\begin{quote}
Say \textquotedblleft{}it will rain\textquotedblright{} and \textquotedblleft{}it rained\textquotedblright{}. (\emph{pāūs paḍel}, \emph{pāūs paḍlā}.) Then \textquotedblleft{}it rained all day\textquotedblright{}. (\emph{divasbhar pāūs paḍat hotā}.)
\end{quote}

\subsection*{You'll want to know: \mr{जर} … \mr{तर}}
\begin{quote}
\textbf{\mr{जर} \mr{उद्या} \mr{पाऊस} \mr{पडला}, \mr{तर} \mr{मी} \mr{घरी} \mr{राहीन}.}
\end{quote}

\begin{quote}
\emph{jar udyā pāūs paḍlā, tar mī gharī rāhīn.} — \textquotedblleft{}If it rains tomorrow, I will stay at home.\textquotedblright{}
\end{quote}

\textbf{\mr{जर}} (\emph{jar}) opens the condition, \textbf{\mr{तर}} (\emph{tar}) opens the result. The condition takes the \textbf{past} form, \textbf{\mr{पडला}}, though it looks ahead; the result is in the future, \textbf{\mr{राहीन}}. \textbf{\mr{जर}} can go: \textbf{\mr{पाऊस} \mr{पडला} \mr{तर} \mr{मी} \mr{घरी} \mr{राहीन}} (\emph{pāūs paḍlā tar mī gharī rāhīn}). \textbf{\mr{तर}} stays.

\subsection*{You'll want to know: if it is}
\begin{quote}
\textbf{\mr{जर} \mr{तुम्हाला} \mr{वेळ} \mr{असेल}, \mr{तर} \mr{या}.}
\end{quote}

\begin{quote}
\emph{jar tumhālā veḷ asel, tar yā.} — \textquotedblleft{}If you have time, come.\textquotedblright{}
\end{quote}

For a state, the condition takes the future of \textquotedblleft{}to be\textquotedblright{}, \textbf{\mr{असेल}}; the result can be a request. \textbf{\mr{या}}, \textquotedblleft{}come\textquotedblright{}, is the polite form of \textbf{\mr{येणे}}.

\subsection*{Your turn}
Say these aloud:

\begin{itemize}
\raggedright
  \item \textquotedblleft{}If it rains tomorrow, I will stay at home\textquotedblright{} — \emph{jar udyā pāūs paḍlā, tar mī gharī rāhīn}
  \item \textquotedblleft{}If Meera comes, we will go\textquotedblright{} — \emph{mīrā ālī tar āmhī jāū}
  \item \textquotedblleft{}If you have time, come\textquotedblright{} — \emph{jar tumhālā veḷ asel, tar yā}
\end{itemize}

\begin{tcolorbox}[breakable,colback=teal!4,colframe=teal!35!black,title={Before you move on}]
\textquotedblleft{}If\textquotedblright{}? (\textbf{\mr{जर}}.) \textquotedblleft{}Then\textquotedblright{}? (\textbf{\mr{तर}}.) Which form does the condition take? (The past: \textbf{\mr{पडला}}.) \textquotedblleft{}If you have time\textquotedblright{}? (\textbf{\mr{जर} \mr{तुम्हाला} \mr{वेळ} \mr{असेल}}.)
\end{tcolorbox}
\section[\texorpdfstring{\mr{पाऊस} \mr{पडल्यास} (pāūs paḍalyās)}{पाऊस पडल्यास (pāūs paḍalyās)}]{\mr{पाऊस} \mr{पडल्यास} (pāūs paḍalyās) — should it happen, and otherwise}

\label{lesson:MR-C286-padlyas}



\begin{quote}
Say \textquotedblleft{}if it rains tomorrow, I will stay at home\textquotedblright{}. (\emph{jar udyā pāūs paḍlā, tar mī gharī rāhīn}.) Then \textquotedblleft{}after reading\textquotedblright{}. (\emph{vācalyāvar}.)
\end{quote}

\subsection*{You'll want to know: -\mr{ल्यास}, 'should it'}
\begin{quote}
\textbf{\mr{पाऊस} \mr{पडल्यास} \mr{आम्ही} \mr{घरी} \mr{राहू}.}
\end{quote}

\begin{quote}
\emph{pāūs paḍalyās āmhī gharī rāhū.} — \textquotedblleft{}Should it rain, we will stay at home.\textquotedblright{}
\end{quote}

The \textbf{-\mr{ल्}-} of \textbf{-\mr{ल्यावर}}, with \textbf{-\mr{ल्यास}}: \textbf{\mr{पडल्यास}} says \textbf{\mr{जर} … \mr{पडला} \mr{तर}} in one word. It is the careful, written way, and you will read it on notices more than you will hear it: \textbf{\mr{उशीर} \mr{झाल्यास}} (\emph{uśīr jhālyās}), \textquotedblleft{}in case of delay\textquotedblright{}.

\subsection*{You'll want to know: \mr{नाहीतर}, 'otherwise'}
\begin{quote}
\textbf{\mr{लवकर} \mr{निघा}, \mr{नाहीतर} \mr{गाडी} \mr{चुकेल}.}
\end{quote}

\begin{quote}
\emph{lavkar nighā, nāhītar gāḍī cukel.} — \textquotedblleft{}Leave early, otherwise you'll miss the train.\textquotedblright{}
\end{quote}

\textbf{\mr{नाहीतर}} (\emph{nāhītar}) is \textbf{\mr{नाही}} and \textbf{\mr{तर}}: \textquotedblleft{}if not, then\textquotedblright{}. \textbf{\mr{निघा}} is the polite \textquotedblleft{}set out\textquotedblright{}; \textbf{\mr{चुकेल}} is the future of \textbf{\mr{चुकणे}}, with the train as its subject, as in \textbf{\mr{माझी} \mr{गाडी} \mr{चुकली}}.

\subsection*{Your turn}
Say these aloud:

\begin{itemize}
\raggedright
  \item \textquotedblleft{}Should it rain, we will stay at home\textquotedblright{} — \emph{pāūs paḍalyās āmhī gharī rāhū}
  \item \textquotedblleft{}Leave early, otherwise you'll miss the train\textquotedblright{} — \emph{lavkar nighā, nāhītar gāḍī cukel}
  \item \textquotedblleft{}Come tomorrow, otherwise I'll get angry\textquotedblright{} — \emph{udyā yā, nāhītar malā rāg yeīl}
\end{itemize}

\begin{tcolorbox}[breakable,colback=teal!4,colframe=teal!35!black,title={Before you move on}]
\textquotedblleft{}Should it rain\textquotedblright{}? (\textbf{\mr{पाऊस} \mr{पडल्यास}}.) Is it said or written more? (Written, on notices.) \textquotedblleft{}Otherwise\textquotedblright{}? (\textbf{\mr{नाहीतर}}.)
\end{tcolorbox}
\section[(review)]{Review — if and otherwise}

\label{lesson:MR-R286-if-and-otherwise}



\begin{quote}
Say \textquotedblleft{}the train is faster than the bus\textquotedblright{} and \textquotedblleft{}I have lived in Pune for three years\textquotedblright{}. (\emph{gāḍī baspekṣā jalad āhe}, \emph{mī tīn varṣāmpāsūn puṇyāt rāhto}.) Then \textquotedblleft{}the biggest\textquotedblright{} and \textquotedblleft{}from morning until evening\textquotedblright{}. (\emph{sarvāt moṭhā}, \emph{sakāḷpāsūn sandhyākāḷparyant}.) Last, \textbf{\mr{ऋतू}} and \textbf{\mr{ज्ञान}}, and their letters \textbf{\mr{ऋ}} and \textbf{\mr{ञ}}.
\end{quote}

\subsection*{Your turn}
Say these aloud:

\begin{itemize}
\raggedright
  \item If it rains tomorrow, I will stay at home.
  \item If you have time, come.
  \item Leave early, otherwise you'll miss the train.
  \item If it rains, you shouldn't go out.
  \item one plan of yours for next week, with \textbf{\mr{जर} … \mr{तर}}
\end{itemize}

\begin{tcolorbox}[breakable,colback=teal!4,colframe=teal!35!black,title={Before you move on}]
\textquotedblleft{}If … then\textquotedblright{}? (\textbf{\mr{जर} … \mr{तर}}.) Which form does the condition take? (The past: \textbf{\mr{पडला}}.) \textquotedblleft{}Should it rain\textquotedblright{}? (\textbf{\mr{पाऊस} \mr{पडल्यास}}.) \textquotedblleft{}Otherwise\textquotedblright{}? (\textbf{\mr{नाहीतर}}.)
\end{tcolorbox}
